\subsection{特征值与特征向量}

定义1. --- 设 E 是交换域 K 上的一个向量空间，u 是 E 的一个自同态。E 的一个元素 x 称为 u 的特征向量，如果存在 \(\alpha \in K\) 使得 \(u \cdot x = \alpha x\)；如果 \(x \neq 0\)，则标量 \(\alpha\) 称为 u 的对应于特征向量 x 的特征值。对每个标量 \(\alpha\)，向量子空间 \(V_{\alpha}\) 由那些 \(x \in E\)（使得 \(u \cdot x = \alpha x\)）组成，它称为 E 的对应于 \(\alpha\) 的特征子空间。

特征值 \(\alpha\) 的几何重数是基数维 \(V_{\alpha}\)。

假设 E 是有限维的。u 的特征值是那些元素 \(\alpha\)（K 中的），它们使得 E 的自同态 \(\alpha \cdot 1 - u\) 不是单射，换言之（III，第 524 页，命题 3），使得 \(\det(\alpha \cdot 1 - u) = 0\)。但由 u 的特征多项式 \(\chi_{u}\) 的定义（III，第 541 页，定义 3），我们有 \(\det(\alpha \cdot 1 - u) = \chi_{u}(\alpha)\)。因此：

命题1. --- 设 E 是有限维的。则 K 的一个元素 \(\alpha\) 是自同态 \(u\) 的特征值，当且仅当它是 \(u\) 的特征多项式的根。

若 L 是域 K 的一个扩域，则 \(\chi_{u}\) 在 L 中的根是自同态 \(1_{L} \otimes u\)（L-向量空间 \(L \otimes_{K} E\) 的自同态）的特征值；它们通常被称为 u 在 L 中的特征值。借用语言习惯，若 u 在 L 的某个代数闭扩张中的所有特征值都属于 L，我们就说 u 的所有特征值都属于 L；这意味着 \(\chi_{u}\) 在 L 中分解为一次因式{[}X{]}。

设 \(U\) 是 \(n\) 阶方阵，其系数在 \(K\) 中。则根据定义，\(U\) 的特征多项式是

\[
\chi_ {U} (\mathbf {X}) = \det \left(\mathbf {X} \cdot I _ {n} - U\right);
\]

U 的特征值（在 K 的扩域 L 中）是多项式 \(\chi_{U}\) 在 L 中的根；它们也是（L 中的）标量 \(\alpha\)：使得线性方程组 \(UX = \alpha X\) 存在非零解，其中 X 是 n 阶列矩阵；满足此方程的列矩阵 X 称为 U 的对应于 \(\alpha\) 的特征向量。

如果 U 是 n 维向量空间的自同态 u 关于基 B 的矩阵，那么 \(\chi_{U} = \chi_{u}\) ，U 的特征值就是 u 的特征值，而 U 的特征向量就是 u 的特征向量关于基 B 的矩阵。

命题2. --- 设 u 是交换域 K 上向量空间 E 的一个自同态；对每个标量 \(\alpha\)，设 \(V_{\alpha}\) 为对应于 \(\alpha\) 的特征子空间。于是子空间 \(V_{\alpha}\) 在 u 下封闭，且这些 \(V_{\alpha}\) 的和是直和。

第一个断言是显然的。根据定义，子空间 \(V_{\alpha}\) 被元素 X − \(\alpha\)（属于 K{[}X{]}）所零化；这些 X − \(\alpha\)（\(\alpha \in K\)）是不可约的且两两不相伴；因此第二个断言由 VII，第 8 页，定理 1 得出。

